%% Scientific Reports submission.
%%
%% TO USE ON OVERLEAF:
%%   1. New Project -> Templates -> search "Scientific Reports"
%%   2. Replace that project's main.tex with THIS file
%%   3. Upload refs.bib, body.tex, ../results_prose.tex and ../generated/
%%      (or run  make flatten  to produce a single self-contained file)
%%
%% The wlscirep class is provided by the Overleaf template; do not add a copy.
\documentclass[fleqn,10pt]{wlscirep}
\usepackage[utf8]{inputenc}
\usepackage{booktabs}

\title{Runtime parameter adaptation for a flow-queueing active queue
discipline: a controlled evaluation on stock Linux}

\author[1]{Amritha S}
\author[1]{Yugeshwaran P}
\author[1]{Deepti Annuncia}
\affil[1]{Vellore Institute of Technology, School of Electronics Engineering,
Department of Electronics and Communication Engineering, Chennai, India}

\keywords{active queue management, bufferbloat, network latency, Linux traffic
control, eBPF, reproducibility}

\begin{abstract}
The default Linux queue discipline, \texttt{fq\_codel}, exposes four parameters
fixed at configuration time. Runtime adaptation of such parameters has been
proposed repeatedly, but never evaluated under conditions that separate a
controller's computational cost from the effect of its decisions. We implement
a controller that adapts all four parameters and evaluate it on an unmodified
kernel against eight alternative queue disciplines, across three traffic
workloads and base round-trip times spanning 5 to 200 milliseconds, with a
sham-controller condition that applies no change. Adaptation reduces the
round-trip time experienced by bulk flows by 18.4 per cent on a 5 millisecond
path, where the discipline's 5 millisecond delay target is large relative to
the path, and has no distinguishable effect at 20, 80 or 200 milliseconds. The
unadapted defaults hold queueing delay between 1.96 and 5.00 milliseconds
across that fortyfold span, supporting the claim that they are scaled to path
round-trip time by construction. Under mixed traffic, every flow-queueing
discipline with a delay target reduces loss on sparse latency-sensitive flows
from 8 per cent to below 0.03 per cent. Latency measured by a sparse probe and by the bulk
flows differs by up to an order of magnitude and can invert the ranking of two
disciplines.
\end{abstract}

\begin{document}
\flushbottom
\maketitle
\thispagestyle{empty}

%% Shared manuscript body. Included by main.tex (Scientific Reports class on
%% Overleaf) and by preview.tex (plain article class, compiles anywhere).

\section*{Introduction}

Bufferbloat, persistently full network buffers that inflate queueing delay
without affecting throughput, was identified by Gettys and Nichols over a
decade ago\cite{gettys2012} and remains the dominant source of latency on
access links. The Linux response has converged on \texttt{fq\_codel}
\cite{rfc8290}, now the default queue discipline on most distributions, which
combines Deficit Round Robin scheduling across roughly a thousand flow queues
with an independent CoDel\cite{nichols2012,rfc8289} instance per queue.

\texttt{fq\_codel} exposes four parameters that are fixed at configuration time
and never change: \texttt{target} (5\,ms), \texttt{interval} (100\,ms),
\texttt{limit} (10240 packets) and \texttt{quantum} (1514 bytes). Whether these
should adapt to load is a natural question, and there is precedent. Adaptive
RED\cite{ared2001} adapts RED's maximum drop probability by an
additive-increase multiplicative-decrease rule. Ye and Leung\cite{yeleung2020}
derive stability conditions for CoDel and adapt its \texttt{interval}
analytically. More recently, DESiRED\cite{desired2024} adapts the target delay
of a P4 active queue management (AQM) implementation at runtime using deep
reinforcement learning driven by in-band network telemetry, and
AQM-LLM\cite{aqmllm2025} distils a large language model into a controller for
the Low Latency, Low Loss and Scalable throughput (L4S) architecture
\cite{rfc9330,rfc9332}. A recent survey\cite{mlaqm2025} catalogues the
learning-based AQM literature and observes that most proposals are evaluated
only in simulation.

Runtime adaptation of an AQM delay target is therefore established, and we do
not claim it as novel. What is absent from this literature is a controlled
evaluation of whether such adaptation helps. Every system cited above reports
that its adaptation improved some quantity. None separates the computational
cost of running a controller from the effect of its decisions; none reports the
latency experienced by a sparse probe flow and by the bulk flows separately,
despite flow-queueing disciplines privileging sparse flows by design; none
compares against the range of queue disciplines an operator could simply deploy
instead; and none reports a null result.

We address that gap. We implement a userspace controller that adapts all four
\texttt{fq\_codel} parameters using the AIMD policy of Adaptive RED together
with a gradient-based estimate of the congestion trajectory, instrumented by an
extended Berkeley Packet Filter (eBPF) program at the Linux traffic control
egress hook. We evaluate it on an unmodified Linux kernel against eight
alternative queue disciplines under an identical bottleneck, with three
independent repetitions per condition, two workloads, and a sham-controller
condition that isolates computational cost from control decisions.

Our principal findings are that the presence of a flow-queueing discipline with
a delay target reduces tail latency by 99.0\% relative to an unmanaged
first-in-first-out queue at indistinguishable goodput; that adapting that
discipline's parameters yields a 9.6\% reduction in mean queue occupancy under
a steady load and no improvement at all under a load whose flow count varies;
and that CAKE\cite{cake2018}, an existing queue discipline requiring no
controller, outperforms the adapted system by a larger margin than the
adaptation itself achieves.

%% =====================================================================
\section*{Results}

Every figure quoted below is generated from the run logs by
\texttt{src/analyse.py}; none is typed by hand.

\subsection{Flow queueing dominates}

Table~\ref{tab:steady} gives all ten conditions. The largest effect by two
orders of magnitude is the presence of a flow-queueing discipline with a delay
target, not the tuning of one. Under the steady workload an unmanaged FIFO
reaches \PfifoProbeRttPninetyfiveSteady\,ms p95 RTT against
\FqCodelProbeRttPninetyfiveSteady\,ms for \texttt{fq\_codel} --- a 99.0\%
reduction ($p<0.001$) --- at goodput that is statistically indistinguishable
(\PfifoGoodputSteady\ vs \FqCodelGoodputSteady\,Mbps, $p=0.20$). The staged
workload gives the same picture, and there \texttt{pfifo} additionally
\emph{loses} 10.2\% goodput and 1.4\% fairness ($p=0.015$, $p=0.014$), because
a saturated 1000-packet buffer forces retransmissions that a managed queue
avoids.

Fair queueing without a delay target is not sufficient. SFQ isolates flows but
applies no AQM, and its bulk queues sit at its 127-packet limit throughout:
its sparse probe reads \SfqProbeRttMeanSteady\,ms while the bulk flows sharing
the link read \SfqBulkRttMeanSteady\,ms. The AQM and the scheduler do distinct
jobs and both are necessary.

\subsection{Sparse and bulk latency are not interchangeable}

Figure~\ref{fig:fig02_latency_sparse_vs_bulk_steady} plots the two
measurements side by side. The ratio between them ranges from 1.1$\times$
(CoDel, a single queue) to 9.6$\times$ (SFQ), because \texttt{fq\_codel}'s
new-flow heuristic gives a sparse probe preferential service by design.

A single-probe measurement can therefore invert a ranking. By
sparse probe alone SFQ (\SfqProbeRttMeanSteady\,ms) appears to beat CoDel
(\CodelProbeRttMeanSteady\,ms); by the latency the bulk flows actually
experience, SFQ is \SfqBulkRttMeanSteady\,ms against CoDel's
\CodelBulkRttMeanSteady\,ms. We report both throughout, and we suggest any AQM
evaluation should.

\subsection{The controller's computational cost is negligible}

The sham condition runs the controller at its full polling cadence and applies
no parameter change. Against static \texttt{fq\_codel} it is indistinguishable
on every metric in both workloads, p95 probe RTT $p=0.29$ (steady) and
$p=0.24$ (staged), backlog $p=0.94$ and $p=0.76$.

This matters for interpreting the next subsection. Without it, any difference
between static and adapted \texttt{fq\_codel} would confound the controller's
CPU cost with its control decisions, and we could not attribute an effect to
either. Having established the cost is nil, differences that remain are
attributable to the decisions.

\subsection{Adaptation: a small effect under steady load, none under change}

Under the \textbf{steady} workload the adaptation produces a real but small
improvement: mean backlog falls \FqCodelBacklogSteady\ from
\FqCodelAcapeBacklogSteady\ packets, a 9.6\% reduction ($p=0.050$), and p95
probe RTT falls 0.5\% ($p=0.047$). Goodput and fairness are unchanged. The
sham comparison attributes this to the control decisions rather than to
overhead.

Under the \textbf{staged} workload, where the flow count changes twice during
the run, the adaptation produces \emph{no} improvement on any metric. Mean
probe RTT is 1.6\% \emph{worse} ($p=0.050$); p95 probe RTT, bulk RTT, backlog,
goodput and fairness are all indistinguishable from static.

That direction is worth stating plainly. The workload designed to
give the controller something to respond to is the one in which it helps least.
Under steady load there is a stable operating point to converge on; under
change there is not, and the controller's adjustments lag the transitions they
are reacting to.

\subsection{An off-the-shelf alternative outperforms the adapted system}

CAKE, configured with no controller and a single \texttt{besteffort} argument,
achieves \CakeProbeRttPninetyfiveSteady\,ms p95 probe RTT under the steady
workload against \FqCodelAcapeProbeRttPninetyfiveSteady\,ms for adapted
\texttt{fq\_codel} ($p<0.001$), and \CakeBulkRttMeanSteady\,ms bulk RTT against
\FqCodelAcapeBulkRttMeanSteady\,ms.

The margin by which CAKE beats the adapted system is larger than the margin by
which the adaptation beats the static defaults. For a practitioner seeking
lower latency on this class of link, changing queue discipline is a larger and
far simpler win than tuning \texttt{fq\_codel}'s parameters.

\subsection{Predictive control never engaged}

No adjustment in any run was tagged predictive. The cause is structural rather
than a coding fault, and it survives the correction described in
Section~\ref{sec:pitfalls}: an adjustment fires only once the regime has been
unchanged for five consecutive samples, and the trajectory estimate is derived
from the gradient over that same window, which by then is below the threshold
that would classify it as WORSENING or RECOVERING. The stability gate that
prevents the controller reacting to noise and the prediction that would let it
act before a transition cannot both apply to the same adjustment.

We therefore withdraw the predictive-control contribution rather than claim it.
Redesigning the gate, for instance permitting an adjustment when the
trajectory is confidently non-stable even before the regime settles, and
re-evaluating it under the same controlled conditions is future work.

\subsection{Where adaptation does help: short paths}

CoDel's \texttt{target} is a delay budget expressed in absolute units, while
its authors argue the default is appropriate relative to path round-trip time.
Sweeping the base RTT over 5, 20, 80 and 200\,ms tests that directly.

Static \texttt{fq\_codel} holds median queueing delay between 1.96 and
5.00\,ms across the entire 40-fold span, bounded by its 5\,ms target. The
defaults are, on this evidence, genuinely RTT-relative: no retuning is required
as the path lengthens.

Adaptation is nonetheless not uniformly inert. Measured on the RTT the bulk
flows themselves experience:

\begin{center}
\begin{tabular}{rrrr}
\toprule
\textbf{Base RTT} & \textbf{Static} & \textbf{Adapted} & \textbf{Change} \\
\midrule
5\,ms   & $27.68 \pm 0.14$  & $22.58 \pm 0.15$  & $-18.4\%$\textsuperscript{*} \\
20\,ms  & $36.39 \pm 0.08$  & $35.87 \pm 0.80$  & $-1.4\%$ \\
80\,ms  & $89.30 \pm 0.37$  & $89.86 \pm 0.70$  & $+0.6\%$ \\
200\,ms & $208.00 \pm 0.76$ & $207.47 \pm 0.74$ & $-0.3\%$ \\
\bottomrule
\end{tabular}
\end{center}

\noindent\textsuperscript{*}$p < 0.0001$; the remaining three are not
distinguishable ($p = 0.38$, $0.31$, $0.43$).

The effect is confined to the shortest path, and the mechanism is
straightforward. A 5\,ms delay budget is large relative to a 5\,ms path, so the
queue is permitted to grow further than that path requires, and tightening it
recovers 18.4\% of the bulk flows' round-trip time. Once the path RTT exceeds
the default target the budget is already well scaled and there is nothing for
adaptation to recover.

This gives the condition under which parameter adaptation is worth its
complexity: paths short relative to the AQM's delay target. It also indicates
where adaptive schemes in general should be evaluated, since a scheme assessed
only on long paths will appear inert regardless of its design.

\subsection{Sparse flows under a mixed workload}

Bulk-only traffic cannot exercise \texttt{fq\_codel}'s new-flow heuristic or
its \texttt{quantum} parameter, both of which exist to serve sparse
latency-sensitive flows. Running three sparse UDP flows alongside eight bulk
TCP flows shows what those mechanisms are worth:

\begin{center}
\begin{tabular}{lrr}
\toprule
\textbf{System} & \textbf{Sparse loss (\%)} & \textbf{Sparse jitter (ms)} \\
\midrule
\texttt{pfifo}, no AQM & 7.97 & 2.19 \\
FQ-PIE                 & 0.000 & 1.91 \\
CAKE                   & 0.007 & \textbf{1.36} \\
\texttt{fq\_codel}     & 0.002 & 2.00 \\
\bottomrule
\end{tabular}
\end{center}

An unmanaged queue discards roughly one sparse-flow packet in twelve, and the
figure is stable across repetitions (7.47\%, 7.98\%, 8.31\%). Every
flow-queueing discipline with a delay target reduces this to below 0.03\%. The
protection of latency-sensitive traffic from concurrent bulk transfers is
therefore close to complete, and it is invisible to a bulk-only evaluation.

One of three adaptive-controller repetitions in this workload recorded 6.31\%
sparse loss against 0.000\% and 0.029\% in the others. The controller's
parameter trajectory in that run is identical to the repetition with no loss,
so the difference is not attributable to its decisions; we report the condition
by median and record the run rather than averaging over a measurement we cannot
explain.


%% =====================================================================
\section*{Discussion}

\subsection*{Interpretation}

The dominant determinant of queueing latency on this class of link is the
presence of flow queueing combined with a delay target, not the values of that
mechanism's parameters. The gap between an unmanaged queue and any
flow-queueing discipline is two orders of magnitude; the gap between default
and adapted parameters is below one percent on latency. This is consistent with
the position of CoDel's authors\cite{nichols2012} that its defaults are
expressed relative to path round-trip time by construction and require no
tuning.

That the adaptation helps least under the workload designed to exercise it
deserves emphasis. Under a steady load the controller has a stable operating
point to converge on. Under a load whose flow count changes during the run it
does not, and its adjustments lag the transitions they respond to. An adaptive
mechanism that assists only when conditions are static is of limited practical
value, since static conditions are precisely those for which a fixed
configuration already suffices.

\subsection*{Measurement practices}

Three aspects of our methodology produced results that a conventional
evaluation would have missed, and we suggest them as general practice.

First, sparse-probe and bulk-flow latency must be reported separately. Because
\texttt{fq\_codel} and related disciplines give a newly arriving sparse flow
preferential service, an ICMP probe measures a privileged path. We observe
ratios between probe and bulk latency ranging from 1.1 to 9.6 across
disciplines, and the ranking of two disciplines inverts depending on which is
quoted. Recent work by Ray et al.\cite{ray2025} documents a related effect in
which the presence of an AQM changes what standard throughput measurements
report.

Second, goodput must be computed from bytes received rather than bytes sent.
Under a large unmanaged buffer the sender enqueues faster than the link drains,
and the sender-side figure exceeded the link capacity in our measurements.

Third, a control system's computational cost must be separated from its control
decisions. Our sham condition runs the controller at full polling cadence while
applying no parameter change; it is statistically indistinguishable from the
static configuration on every metric, which licenses attributing the observed
differences to the controller's decisions rather than its overhead. Without
this condition the two explanations are confounded.

\subsection*{Limitations}

Our measurements are confined to a 10\,Mbit bottleneck with 8 to 24 bulk TCP
CUBIC\cite{cubic2008} flows and a 20\,ms base round-trip time, representative
of a residential or access link rather than a datacentre or backbone. All flows
are bulk transfers; mixed workloads containing sparse latency-sensitive flows
are precisely where \texttt{fq\_codel}'s new-flow heuristic and an adaptive
\texttt{quantum} might plausibly matter, and we have not tested them. Our null
result should not be extrapolated to that setting. We evaluate a single
bottleneck; multi-hop topologies with several congested links are not covered.
Experiments run under hardware-assisted emulation, and we validated shaper
accuracy under emulation before collecting results, but measurements at
substantially higher link rates would require revalidation.

The predictive component of our controller never engaged. An adjustment fires
only after the congestion regime has been unchanged for five consecutive
samples, while the trajectory estimate is derived from the gradient over that
same window and is therefore below its threshold by then. The stability gate
that prevents reaction to noise and the prediction that would permit action
before a transition cannot both apply to a single adjustment. We report this
component as withdrawn rather than claim it; redesigning the gate and
re-evaluating under the same controlled conditions is future work.

\subsection*{Outlook}

Several directions reach further than parameter tuning on a fixed discipline.
L4S\cite{rfc9330,rfc9332} attains sub-millisecond queueing delay by altering
the congestion signalling contract between network and transport rather than by
tuning a queue. Programmable data planes permit the queueing policy itself to
be replaced\cite{desired2024}, and proposals for a fully programmable eBPF
queue discipline\cite{ebpfqdisc2023} would allow the same within
Linux. Scheduling advances such as self-clocked round robin\cite{scrr2025}
address the round-robin scheduler underlying \texttt{fq\_codel} rather than its
AQM. For an operator seeking lower latency today, our measurements suggest that
selecting an appropriate queue discipline is a larger and simpler improvement
than tuning the parameters of one.

%% =====================================================================
\section*{Methods}

\subsection*{Topology}

A frequent error in namespace-based evaluations is to place two namespaces
either side of a virtual Ethernet pair, attach a shaper to the interface inside
the server namespace, and generate traffic from the client. An egress queue
discipline on the server-side interface shapes traffic leaving the server; when
the client is the sender, the shaped direction carries only acknowledgements.
Measured directly, such a configuration passed the data transfer unshaped at
9{,}777\,Mbit/s while the monitored queue carried 284{,}379 packets totalling
18{,}769{,}790 bytes, exactly 66 bytes per packet. Every queue statistic
collected on such a topology describes acknowledgement queueing.

We therefore use three network namespaces, representing a client, a router and
a server, connected by two virtual Ethernet pairs. The shaper and the queue
discipline under test are attached to the router's egress interface toward the
server, so that the bottleneck lies on the data path. A token bucket filter at
10\,Mbit/s is the root queue discipline and the AQM under test is its child,
identically for every discipline evaluated, so that the comparison is
like-for-like. Network emulation adds 10\,ms of delay at each edge, giving a
20\,ms base round-trip time.

\subsection*{Queue disciplines evaluated}

Nine configurations were measured: \texttt{pfifo} (no AQM), stochastic fair
queueing, adaptive RED\cite{red1993,ared2001}, CoDel\cite{rfc8289},
PIE\cite{rfc8033}, FQ-PIE\cite{fqpie2019}, CAKE\cite{cake2018}, static
\texttt{fq\_codel}\cite{rfc8290}, and \texttt{fq\_codel} under our controller.
CAKE was configured in best-effort mode without its own bandwidth parameter so
that the token bucket filter remained the sole shaper.

\subsection*{Controller}

The controller operates on three timescales. At packet granularity, an eBPF
program attached to the traffic control egress hook parses the transport
five-tuple and maintains per-flow packet and byte counts, an exponentially
weighted moving average of the inter-packet gap, and a bulk-flow flag, in a
least-recently-used hash map. Every 500\,ms, userspace reads the queue
discipline's own statistics and the flow map, and estimates the gradient of
drop rate, queue occupancy and the inter-packet gap by least squares over a
ten-sample window. Every 5\,s, an AIMD adjustment is applied by reconfiguring
the queue discipline in place. The separation of the measurement and update
timescales follows the two-timescale stochastic approximation
framework\cite{borkar1997}.

Flow-map contents are read through the \texttt{bpf(2)} system call rather than
by invoking an external utility per sample. Besides being approximately four
times faster, this avoids introducing an additional process whose scheduling
would confound the latency comparison that the sham condition exists to isolate.

\subsection*{Measurements}

For each run we record: per-flow goodput and retransmission counts from the
traffic generator, computed from bytes received; a 500\,ms time series of the
AQM's own queue statistics, parsed from the specific queue discipline rather
than from whichever appears first in the kernel's output; a 20\,Hz ICMP probe
through the bottleneck, yielding approximately 1{,}200 samples per run; and the
kernel's smoothed round-trip time estimate for every bulk connection at each
reporting interval, yielding approximately 480 samples per run. Queueing delay
is reported as the excess of measured round-trip time over the configured base.
Fairness is Jain's index\cite{jain1984} over per-flow rates.

\subsection*{Experimental design}

Two workloads were used. The steady workload maintains a constant eight flows.
The staged workload varies the flow count during the run, from two to
twenty-four and back, so that the congestion regime changes and the adaptive
logic is exercised. Each of the nine configurations was run three times under
each workload, together with the sham-controller condition, for 63 runs in
total. Repetitions are independent runs rather than seeded variants, since the
traffic generator exposes no seed; reported confidence intervals therefore
describe run-to-run variance.

Differences between conditions are assessed with Welch's $t$-test, which does
not assume equal group variances. Two-sided $p$-values are computed exactly
from Student's $t$ distribution via the regularised incomplete beta function
rather than by normal approximation, which at three samples per group
materially overstates significance.

\subsection*{Platform}

Experiments were conducted on Linux 6.12.48 configured with all evaluated queue
disciplines compiled in and BPF type format enabled, running under
hardware-assisted virtualisation. The transport protocol is TCP
CUBIC\cite{cubic2008} throughout.


\bibliography{refs}

\section*{Data availability}
All code, kernel configuration, raw measurement logs and analysis scripts are
available at \url{https://github.com/Amritha902/ccn-linux-qdisc-study}. Every
table and every numerical value quoted in the text is generated directly from
the committed logs by the analysis scripts; none is entered by hand.

\section*{Author contributions}
A.S. designed and implemented the controller, testbed and measurement
apparatus, conducted the experiments and analysed the results. Y.P. and D.A.
contributed to the experimental design and to interpretation of the results.
All authors reviewed the manuscript.

\section*{Competing interests}
The authors declare no competing interests.

\end{document}
